\documentclass[letterpaper]{article} % DO NOT CHANGE THIS
\usepackage[submission]{aaai2027}  % DO NOT CHANGE THIS
% The serif, sans-serif, and monospaced fonts are loaded automatically by
% aaai2027.sty (newtxtext, helvet, courier). DO NOT add \usepackage{times},
% \usepackage{helvet}, \usepackage{courier}, or any other font package.
\usepackage[hyphens]{url}  % DO NOT CHANGE THIS
\usepackage{graphicx} % DO NOT CHANGE THIS
\urlstyle{rm} % DO NOT CHANGE THIS
\def\UrlFont{\rm}  % DO NOT CHANGE THIS
\usepackage{natbib}  % DO NOT CHANGE THIS AND DO NOT ADD ANY OPTIONS TO IT
\usepackage{caption} % DO NOT CHANGE THIS AND DO NOT ADD ANY OPTIONS TO IT
\frenchspacing  % DO NOT CHANGE THIS
%
% These are recommended to typeset algorithms but not required. See the subsubsection on algorithms. Remove them if you don't have algorithms in your paper.
\usepackage{algorithm}
\usepackage{algorithmic}

%
% These are recommended to typeset listings but not required. See the subsubsection on listing. Remove this block if you don't have listings in your paper.
\usepackage{newfloat}
\usepackage{listings}
\DeclareCaptionStyle{ruled}{labelfont=normalfont,labelsep=colon,strut=off} % DO NOT CHANGE THIS
\lstset{%
	basicstyle={\footnotesize\ttfamily},% footnotesize acceptable for monospace
	numbers=left,numberstyle=\footnotesize,xleftmargin=2em,% show line numbers, remove this entire line if you don't want the numbers.
	aboveskip=0pt,belowskip=0pt,%
	showstringspaces=false,tabsize=2,breaklines=true}
\floatstyle{ruled}
\newfloat{listing}{tb}{lst}{}
\floatname{listing}{Listing}

%
% Recommended for better-looking tables
\usepackage{booktabs}

\usepackage{tikz}
\usetikzlibrary{arrows.meta,positioning}
\newcommand{\op}{\diamond}

\usepackage{amsmath}

%
% Keep the \pdfinfo as shown here. There's no need
% for you to add the /Title and /Author tags.
\pdfinfo{
/TemplateVersion (2027.1)
}

% DISALLOWED PACKAGES
% \usepackage{authblk} -- This package is specifically forbidden
% \usepackage{balance} -- This package is specifically forbidden
% \usepackage{CJK} -- This package is specifically forbidden
% \usepackage{float} -- This package is specifically forbidden
% \usepackage{flushend} -- This package is specifically forbidden
% \usepackage{fullpage} -- This package is specifically forbidden
% \usepackage{geometry} -- This package is specifically forbidden
% \usepackage{hyperref} -- This package is specifically forbidden
% \usepackage{navigator} -- This package is specifically forbidden
% (or any other package that embeds links such as navigator or hyperref)
% \usepackage{indentfirst} -- This package is specifically forbidden
% \usepackage{layout} -- This package is specifically forbidden
% \usepackage{multicol} -- This package is specifically forbidden
% \usepackage{nameref} -- This package is specifically forbidden
% \usepackage{savetrees} -- This package is specifically forbidden
% \usepackage{setspace} -- This package is specifically forbidden
% \usepackage{stfloats} -- This package is specifically forbidden
% \usepackage{tabu} -- This package is specifically forbidden
% \usepackage{titlesec} -- This package is specifically forbidden
% \usepackage{tocbibind} -- This package is specifically forbidden
% \usepackage{ulem} -- This package is specifically forbidden
% \usepackage{wrapfig} -- This package is specifically forbidden
% DISALLOWED COMMANDS
% \nocopyright -- Your paper will not be published if you use this command
% \addtolength -- This command may not be used
% \balance -- This command may not be used
% \baselinestretch -- Your paper will not be published if you use this command
% \clearpage -- No page breaks of any kind may be used for the final version of your paper
% \columnsep -- This command may not be used
% \newpage -- No page breaks of any kind may be used for the final version of your paper
% \pagebreak -- No page breaks of any kind may be used for the final version of your paper
% \pagestyle -- This command may not be used
% \tiny -- This is not an acceptable font size.
% \vspace{- -- No negative value may be used in proximity of a caption, figure, table, section, subsection, subsubsection, or reference
% \vskip{- -- No negative value may be used to alter spacing above or below a caption, figure, table, section, subsection, subsubsection, or reference

\setcounter{secnumdepth}{0} %May be changed to 1 or 2 if section numbers are desired.

% The file aaai2027.sty is the style file for AAAI Press
% proceedings, working notes, and technical reports.
%

% Title

% Your title must be in mixed case, not sentence case.
% That means all verbs (including short verbs like be, is, using,and go),
% nouns, adverbs, adjectives should be capitalized, including both words in hyphenated terms, while
% articles, conjunctions, and prepositions are lower case unless they
% directly follow a colon or long dash
\title{ALPS: A Renewable, Verifiable Benchmark for \\ Mathematical Construction Beyond Automated Reach}
\author{
    %Authors
    % All authors must be in the same font size and format.
    Written by AAAI Press Staff\textsuperscript{\rm 1}\thanks{With help from the AAAI Publications Committee.}\\
    AAAI Style Contributions by Peter Patel Schneider,
    Sunil Issar,\\
    J. Scott Penberthy,
    George Ferguson,
    Hans Guesgen,
    Francisco Cruz\equalcontrib\corresponding,
    Marc Pujol-Gonzalez\equalcontrib\corresponding
}
\affiliations{
    %Afiliations
    \textsuperscript{\rm 1}Association for the Advancement of Artificial Intelligence\\
    % If you have multiple authors and multiple affiliations
    % use superscripts in text and roman font to identify them.
    % For example,

    % Sunil Issar\textsuperscript{\rm 2},
    % J. Scott Penberthy\textsuperscript{\rm 3},
    % George Ferguson\textsuperscript{\rm 4}\corresponding,
    % Hans Guesgen\textsuperscript{\rm 5}
    % Note that the comma should be placed after the superscript

    1101 Pennsylvania Ave, NW Suite 300\\
    Washington, DC 20004 USA\\
    % email address must be in roman text type, not monospace or sans serif
    proceedings-questions@aaai.org
%
% See more examples next
}

%Example, Single Author, ->> remove \iffalse,\fi and place them surrounding AAAI title to use it
\iffalse
\title{My Publication Title --- Single Author}
\author {
    Author Name
}
\affiliations{
    Affiliation\\
    Affiliation Line 2\\
    name@example.com
}
\fi

\iffalse
%Example, Multiple Authors, ->> remove \iffalse,\fi and place them surrounding AAAI title to use it
\title{My Publication Title --- Multiple Authors}
\author {
    % Authors
    First Author Name\textsuperscript{\rm 1,\rm 2}\equalcontrib,
    Second Author Name\textsuperscript{\rm 2}\equalcontrib,
    Third Author Name\textsuperscript{\rm 1}\corresponding
}
\affiliations {
    % Affiliations
    \textsuperscript{\rm 1}Affiliation 1\\
    \textsuperscript{\rm 2}Affiliation 2\\
    firstAuthor@affiliation1.com, secondAuthor@affilation2.com, thirdAuthor@affiliation1.com
}
\fi


\begin{document}

\maketitle

\begin{abstract}

\end{abstract}

% Uncomment the following to link to your code, datasets, an extended version or similar.
% You must keep this block between (not within) the abstract and the main body of the paper.
% Make sure that you do not de-anonymize yourself with these links.
% \begin{links}
%     \link{Code}{https://aaai.org/example/code}
%     \link{Datasets}{https://aaai.org/example/datasets}
%     \link{Extended version}{https://aaai.org/example/extended-version}
% \end{links}

\section{Introduction}
% Why is this benchmark useful? - AI for science benchmarks should be verifiable, non-vacuous, contamination-free, and method-bound


As large language models (LLMs) are applied across mathematics and the sciences, AI-for-science has grown increasingly interested in evaluating LLM reasoning and creativity. Recently, automated mathematical reasoning has become one of the most visible frontiers within AI-for-science, with systems now reaching gold-medal performance at the International Mathematical Olympiad \citep{deepmind2025imo}. Two major recent developments shift this attention specifically towards generating proofs. In 2026, a general-purpose reasoning model resolved the roughly eighty-year-old Erdős unit-distance problem, a long-open question in discrete geometry \citep{openai2026unitdistance}. The SAIR Foundation's Mathematics Distillation Challenge \citep{tao2026distillation}, organized by Terence Tao and Damek Davis, asks competitors to design LLM systems that can construct proofs over the algebraic questions of the Equational Theories Project \citep{bolan2025equational}. The common thread is a standard that the rest of AI-for-science rarely meets: a creative answer only counts once it has been proven correct. 

\input{Figures/fig_introduction}

Verifiability of this kind is just one component necessary for a rigorous benchmark. As LLMs begin to tackle scientific work - formalizing and proving theorems \citep{yang2023leandojo, xin2024deepseek}, proposing conjectures \citep{romera2024mathematical, gottweis2026accelerating, lu2024ai}, and designing molecules and materials \citep{jumper2021highly, merchant2023scaling} - the essential question becomes whether what they produce is genuinely new. However, current benchmarks fail to sufficiently measure this capability. Judging LLM outputs using human experts or through physical experiment is slow, costly, and subjective. Scoring against reference answers provides an objective measurement, but such answers often contaminate LLM training data \citep{xu2024benchmark, chen2025recent}. This makes it difficult to discern whether a high score reflects recall rather than discovery. What is missing is a benchmark that scores novel work objectively, on problems no LLM could have seen before. 


Problems of this kind exist within equational theory \citep{birkhoff1935structure}. This area of mathematics studies magmas: sets with a single binary operation $\op$ that combines two elements into a third \citep{burris1981course}. Its laws are universally quantified equations that an operation may or may not obey, such as commutativity $x \op y = y \op x$ \citep{burris1981course, baader1998term}. A magma is a model of a law when the equation holds under every assignment of its elements \citep{burris1981course}. The one-element magma satisfies every law and is always a model. If there exists only one element, all variables must always be equivalent to each other. Therefore, this model is never informative, and is therefore named the trivial magma \citep{burris1981course}. What carries content is whether a law has a nontrivial model, one with at least two elements. On a finite carrier, such a model can simply be found by search, since only finitely many operations exist to test. However, some laws have no nontrivial finite model at all. Their only nontrivial models are infinite, and no search through finite candidates can reach an infinite structure \citep{kisielewicz1997austin, bolan2025equational}. Here, the questions stop being a matter of computation and become a matter of reasoning about the magma's structure as a whole.

We build the benchmark from the laws that require this reasoning: those with no nontrivial finite model, a property that we verify in advance. With finite models excluded, each law falls into one of two cases - either it forces every model to be trivial, entailing $x = y$, or it has a nontrivial model, which must then be infinite. Laws of the second kind, an infinite model but no finite one, are the Austin laws \citep{kisielewicz1997austin}. 

Austin laws sit where AI-for-scientists and mathematical provers are converging. The recent Erdős unit-distance result was settled by a reasoning model through the construction of an infinite object, the kind of scientific leap that marks the frontier of what today's language models can do \citep{openai2026unitdistance}. Computational mathematical provers approach this same frontier from the other side, with recent progress demonstrating new methods of identifying proofs for infinite models \citep{janota2026case}. Austin laws are this challenge in a controlled, renewable form: each is an infinite-construction problem with a definite answer, and fresh ones can be generated without limit. It is that combination, a frontier problem that is nonetheless checkable and endless in supply, that makes them something a benchmark can be built upon. 

We call this benchmark Austin-Law Proof-Synthesis (ALPS). Each problem is one such law, and solving one means providing a proof showing which of the two cases holds true. A solver must either establish that the law is Austin, by constructing an infinite model and showing it satisfies the law, or prove that the law entails $x = y$, meaning that all models for this law must be trivial. 

Three properties set ALPS apart, shown as the three overlapping ovals in Figure \ref{fig:benchmark_gap}, where each benchmark family lies within the region matching the properties it has. ALPS is verifiable: each answer is settled by a certificate a machine can check, so grading is human-free and provably correct. It is renewable: fresh instances can be produced without limit, so evaluation always has problems no model has seen. And it demands creative construction: a nontrivial model exists only among infinite structures and cannot be reached with a general procedure, so solutions must be invented for each law. Every prior benchmark family has at most two such properties. Synthetic puzzles such as random SAT \citep{selman1996generating} and the Equational Theories Project \citep{bolan2025equational} are verifiable and renewable, but their problems are solved by a general procedure given enough compute, not by construction. Static corpora such as MATH \citep{hendrycks2021measuring}, coding benchmarks \citep{chen2021evaluating}, and miniF2F \citep{zheng2021minif2f}, ask for constructed, checkable solutions. However, these come from fixed sets that eventually contaminate without a clear way of generating high-quality new problems. Open scientific discovery \citep{lu2024ai, gottweis2026accelerating} demands construction and never exhausts its supply, but its answers are judged subjectively. ALPS alone lies where all three meet. 




What ALPS requires of the solver is a capability that reaches well beyond equational theory. Resolving a law requires proposing a completely novel structure, then arguing from the law alone that it satisfies a specification over infinitely many cases it can never check directly. This mirrors mathematical discovery and the building of scientific theories: the answer can never be searched for or recalled, only hypothesized and then justified. What sets ALPS apart is that here the justification is a formal proof, so the creative leap is isolated and made objectively measurable. These problems are exceedingly difficult, but that is the intent: progress on them is direct evidence of the ability to construct and verify genuinely new structure - the capability that separates discovery from recall, and the one current benchmarks cannot measure.


To evaluate the true difficulty of this benchmark, we construct a portfolio of the strongest automated provers \citep{kovacs2013first, schulz2019faster, smallbone2021twee}, finding that it resolves only a small part of ALPS. A large and growing set of problems resists it entirely, and on these it constructs no model at all, at any compute budget, even though every instance is certified to have a determinate answer. This difficulty is bound by method, not computational budget. In full, our contributions are:
\begin{itemize}
\item a two-sided construction task on Austin laws, with every instance certified to have a determinate answer;
\item a corpus of such laws and a public generator that mints fresh instances without limit;
\item an automated judge that certifies submissions through independent channels - a Lean proof\citep{moura2021lean} for triviality, and an automated-prover certificate \citep{kovacs2013first} for constructed models - with no human in the loop; and
\item a strong automated baseline that establishes the method-bound hard tier and, along the way, settles an order-5 case left open by the Equational Theories Project (ETP) \citep{bolan2025equational}.
\end{itemize}


\section{Related Works}

\subsection{AI-for-Science Benchmarks}

Existing benchmark families each hold two of the three properties ALPS targets and fall short on the third. Static problem sets are fixed collections spanning grade-school and competition mathematics \citep{cobbe2021training, hendrycks2021measuring}, graduate science \citep{rein2023gpqa}, coding \citep{chen2021evaluating}, and formal proof construction \citep{zheng2021minif2f, azerbayev2023proofnet, tsoukalas2024putnambench}. These are scored objectively and require reasoning through each problem rather than recalling a memorized algorithm. However, because their problems are fixed, these benchmarks inevitably suffer from contamination, as their evaluation set will eventually enter the training data of the language models they are meant to test, rendering further evaluation on this benchmark unreliable \citep{xu2024benchmark, chen2025recent}. Some benchmarks attempt to delay contamination by holding problems secret \citep{glazer2024frontiermath} or releasing new problems over time \citep{white2025livebench, jain2025livecodebench}, but will always be limited by the rate at which humans can author new problems.

Synthetic puzzles avoid this limit by generating problems from a fixed pattern or rule so that fresh instances can be produced without end, such as random SAT formulas \citep{selman1996generating}, graph-generated reasoning tasks \citep{zhu2024dyval}, Sudoku puzzles \citep{seely2025sudoku}, or the equational laws of the Equational Theories Project \citep{bolan2025equational}. However, the same regularity that makes these problems generable also makes them solvable by a known procedure. Therefore, these benchmarks often test the LLM's ability to execute a known method rather than to truly construct a creative solution. 

Open scientific discovery benchmarks evaluate an LLM's ability to construct novel, useful scientific advancements such as hypotheses, conjectures, designs, and mathematical structures \citep{lu2024ai, gottweis2026accelerating, alkan2025survey, romera2024mathematical, merchant2023scaling, jumper2021highly}. These pose a challenging task, requiring the LLM to invent a new solution rather than retrieve or execute a known procedure. However, their results are validated by human expertise or physical experiment, making evaluation subjective and slow. ALPS uniquely combines the strengths of each benchmark family: objective scoring, endless renewal, and a demand for genuine construction.


\subsection{Theorem Proving and Verification}
ALPS scores answers using machine verification, without a human in the loop, made possible by theorem provers. Lean \citep{moura2021lean} is a functional programming language equipped with a proof assistant whose kernel can automatically check the validity of proofs, such that a certificate is guaranteed to be correct. Each proposition, including the statement that a magma satisfies a law, is encoded as a type within the language, and a proof of it is a term inhabiting that type, exactly as a value inhabits a type in an ordinary programming language. Checking a proof is resolved using type checking: the kernel confirms that submitted terms have the claimed type, similar to a compiler checking that a function's body matches its signature. A statement quantified over infinitely many inputs, such as a law holding under every assignment, becomes a function type. Within Lean, its proof is a function returning a proof for an arbitrary argument, and the kernel type-checks that function once, symbolically. A certificate that type-checks successfully is therefore correct by construction and automatically validated. 

Unlike checking them, finding proofs is difficult. Many works investigate automating this search. Classical automated theorem provers such as Vampire and E \citep{kovacs2013first, schulz2019faster} search for first-order proofs by superposition calculus \citep{bachmair1994rewrite}, a strategy that repeatedly derives logical consequences of the input until it reaches a contradiction or no new consequences remain. More recently, language models have been trained to construct proofs in systems like Lean \citep{yang2023leandojo, xin2024deepseek, hubert2026olympiad}, measured on formal proof benchmarks \citep{zheng2021minif2f}. 


\subsection{Equational Reasoning}
Equational theory studies the foundation of algebra - the identities that operations satisfy and the entailments among them. The Equational Theories Project \citep{bolan2025equational} studies this at scale, enumerating the implications among thousands of small laws and settling them automatically using a pipeline of automated provers and finite-model finders. This work singles out the laws whose only nontrivial models are infinite, the Austin laws \citep{kisielewicz1997austin}, as the hardest to resolve, since neither a searchable finite model nor a short derivation settles them. 

Several automated methods are tailored to equational laws. Paradox and Mace4 \citep{claessen2003new, mccune2003mace4} are designed to find finite models when they exist, while Infinox \citep{claessen2011automated} instead proves that no finite model exists. Twee \citep{smallbone2021twee} implements unfailing Knuth-Bendix completion \citep{knuth1970simple, bachmair1989completion}, orienting equations into rewrite systems to solve equational problems. When a saturation from Twee or a superposition prover halts without reaching a contradiction, the clauses it has accumulated already describe a model. \citet{janota2026case} recently studied this for equational theories, extracting an explicit model from the saturation in a case study. 


\section{The Austin-Law Proof-Synthesis Benchmark}

\input{Figures/fig_methodology_filters}

\input{Figures/fig_methodology_extension}

ALPS is built in three parts, displayed in Figures \ref{fig:method_filters} and \ref{fig:method_extension}: a two-sided task posed over laws certified to have no nontrivial finite model, a corpus construction that mints such laws at any order by extending known Austin laws and screening the results, and an automated judge that verifies answers to either side of the task with a proof certificate. We describe each component below.

\subsection{Problem Formulation}
% Problem Formulation
% Definitions
% Certificate
% Dichotomy
% Asymmetry and non-vacuity
% Two-sided task


A magma is a set $M$, called the carrier, together with a single binary operation ${\op}\colon M \times M \to M$ and no further axioms. An equational law is a universally quantified identity between two terms built from variables and $\op$. Its order is the number of times the operation appears, so the commutative law $x \op y = y \op x$ has order two. The order of a law is the number of times the operation appears in it. A magma \textit{satisfies} a law when the identity holds under every assignment of carrier elements to its variables, and is then called a \textit{model} of the law.
Every ALPS instance is a law of the form $L\colon x = T[x, y, z, \dots]$, a single variable equated with a term $T$ of order five or more. The trivial magma contains only one element, and is thus a model of every law, carrying no information. We call a model nontrivial when its carrier holds at least two distinct elements. The question ALPS poses for each $L$ is whether a nontrivial model exists at all: either $L$ entails $x = y$, collapsing every model to a single value, or some nontrivial magma satisfies it.


A nontrivial model, when one exists, may be finite or infinite. The finite case is a matter of search: enumerating operations on carriers of increasing size must eventually find it, requiring none of the law-specific creativity the benchmark aims to measure. We therefore admit a law into the benchmark only when it carries a machine-checked proof that no nontrivial finite model exists, calling such laws \textit{admissible}. Admissibility forces a dichotomy - every law either entails $x=y$ or it does not, so each instance has a determinate answer. For an admissible law, the two cases are exactly

\begin{equation}
    \label{eq:dichotomy}
    \underbrace{L \models x = y}_{\text{trivial}}
    \quad\text{or}\quad
    \underbrace{L \text{ has an infinite nontrivial model}}_{\text{Austin}}
\end{equation}
Laws of the second kind are the Austin laws \citep{kisielewicz1997austin}.

The asymmetry between the two sides of the task is essential to what the benchmark measures. Whether $L$ entails $x = y$ is a first-order consequence, so a complete proof search will confirm it whenever it holds. The trivial side is, in principle, settled by deduction. The Austin case admits no such guarantee. A search that fails to derive $x = y$, however long it runs, is not a model. Because equational entailment is undecidable in general \citep{baader1998term}, no uniform procedure can be expected to close the gap. Establishing that a law is Austin means inventing an infinite magma tailored to the individual law to satisfy it. This synthesis, which deductive search does not supply, is the capability that ALPS isolates. 


\subsection{Corpus Construction}
% Construction
% Generation, distinctness, contamination free
\textbf{Generation by extension.} Admissible laws are too rare for random sampling to produce in quantity, and the few it does are often redundant. We instead grow the corpus outward from laws already known to be Austin (Figure \ref{fig:method_extension}): a seed law $x = T$ is extended by selecting any variable $v$ in $T$ and replacing one occurrence with $v \op w$ or $w \op v$ for some other variable $w$, yielding higher-order candidates. Because a candidate inherits most of the seed's term structure, it often inherits the mechanism that prevents a finite model, increasing the likelihood of finding new admissible laws relative to random sampling. The construction is recursive - laws confirmed Austin at order $n$ seed the round at order $n+1$ - so the corpus has no terminal order. 

\textbf{Screening.} Each candidate then passes the screening pipeline shown in Figure \ref{fig:method_filters}. A finite-model filter searches carriers of increasing size, discarding every candidate it can satisfy nontrivially. The survivors face the admissibility prover, a first-order query whose success is a proof that no nontrivial finite model exists: when some subterm of $T$ contains every occurrence of $x$, that subterm acts as a left inverse, making the map $x \mapsto T$ injective. Finite carriers imply surjectivity from injective maps, and asserting that surjectivity lets the prover derive $x = y$. Admission requires this proof, and no such proof exists for a law with a finite model. The cost is incompleteness: when injectivity cannot be proven within a candidate, nothing is certified. Such candidates are discarded, even though some may be admissible. The admissible pool that remains is split by \eqref{eq:dichotomy} into laws proved trivial, laws proven Austin by a model the screening provers construct outright, and the residual that the baseline of the next section attacks. 

\textbf{Deduplication.} Because every law descends from a small seed set, extensions risk equivalence with their seed or with a sibling. For the laws proved Austin, the models found during screening can prove distinctness at no prover cost. A magma in which one law holds and the other fails is already a certificate that the two are distinct. Of the 34,191 pairs across the 262 laws proved Austin, 33,936 are separated by such a model, while 255 are proved equivalent by the prover, with none left undecided. Closing the 255 equivalent pairs under transitivity merges away 67 laws, so the 262 laws collapse to 195 classes. We report corpus size in classes rather than laws so the count only reflects distinct problems. Across class boundaries, the model built for one law satisfies 

\subsection{Answer Verification}
% Answer Verification
% lean, generated goals, safeguards

A solved ALPS instance is a verified mathematical fact, not a predicted label. The two sides of the task call for different verification processes, so the judge issues a certificate through the appropriate channel: a proof checked by the Lean kernel for the trivial side, and an automated-prover certificate for the constructed model. Both channels are fully automatic, and a submission is accepted only if its check succeeds. 

\textbf{The trivial channel.} For each law the judge generates the proposition
\begin{equation}
\label{eq:trivialgoal}
\textsc{TrivialGoal} \equiv \forall M\ \forall {\op}\colon M \times M \to M,\;
\mathrm{Law}(\op) \rightarrow \forall a, b : M,\ a = b,
\end{equation}
where $\mathrm{Law}$ asserts that $\op$ satisfies $L$, and a submission is a Lean proof of it. The judge generates the statement, and the submitted proof must elaborate at exactly that type, so a proof the kernel accepts is correct by construction \citep{moura2021lean}.

\textbf{The construction channel.} A solver that believes $L$ is Austin does not write a Lean proof about an infinite object. It submits the model itself, as a finite presentation $E$ - a set of equations over $\op$ that pins down the intended structure. The judge then runs two independent prover queries
\citep{kovacs2013first}:
\begin{equation}
    \label{eq:construction}
    \text{(a)}\;\; E \vdash L
    \qquad\qquad
    \text{(b)}\;\; E \cup \{\,a \neq b\,\} \text{ is satisfiable,}
\end{equation}
certifying the submission only when both succeed. Query (a) is a refutation proof that every model of $E$ satisfies $L$; query (b) is a terminating saturation establishing that $E$ has a model with two distinct elements. Together they exhibit a nontrivial model of $L$, which admissibility forces to be infinite, so $L$ is Austin.

A single Lean judge for both sides is not available, for a reason that strengthens the benchmark's premise. If an admissible law had a model whose operation is a polynomial or affine formula over $\mathbb{Z}$, $\mathbb{Z}^k$, or $\mathbb{Z}/n$, the identity would survive reduction modulo two and yield a nontrivial finite model, contradicting admissibility. No Austin model is an arithmetic formula, therefore, the formalized model families are provably empty here. Term models built from saturations could validate such a proof in Lean, but they await a formalization of ground confluence \citep{janota2026case}.

To make the solution for each law $L$ checkable, we generate two
propositions a solver may prove. Constructing a model means proving
\[
  \begin{aligned}
    \textsc{AustinGoal} \;\equiv\;{}& \exists\, M\; \exists\, {\op}\colon M \times M \to M,\\
    & \bigl(\exists\, a, b : M,\ a \neq b\bigr) \wedge \mathrm{Law}(\op).
  \end{aligned}
\]
 an existential over a carrier $M$, an operation on it, and a witness that two of its elements differ, where $\mathrm{Law}(\op)$ represents the assertion that an operation $\op$ satisfies $L$. Proving the law trivial means proving
\[
\begin{aligned}
      \textsc{TrivialGoal} \;\equiv\; \forall\, M\; \forall\, {\op}\colon M \times M \to M,\; \\
  \mathrm{Law}(\op) \rightarrow \forall\, a, b : M,\ a = b.
\end{aligned}
\]
The nontriviality clause $\exists a, b.\ a \neq b$ is what separates a genuine model from the one-element magma. Together with the admissibility certificate, it forces that
model to be infinite. Notably, a nontrivial model need not be infinite, whereas Austin laws are defined by their infinite models. However, admissibility has already excluded every nontrivial finite model, so any magma that satisfies $\mathrm{Law}(\op)$ with two distinct elements is necessarily infinite, and a proof of $\textsc{AustinGoal}$ therefore establishes that $L$ is an Austin law. Requiring only nontriviality keeps the goal tractable, since the solver never has to prove infinitude inside Lean.


\input{Figures/table_counts}

\input{Figures/table_distinctness}

\section{Experiments}

\subsection{Baseline Portfolio}
%Baseline - 8 configs
% Reproducibility
We establish the difficulty of the benchmark against a strong portfolio of automated off-the-shelf provers. Because the task is two-sided, the portfolio must be as well. Proving a law trivial requires a proof search, while finding a nontrivial model is completed through saturation. The portfolio spans eight configurations: Vampire~5.0.1 in five modes (proof search and saturation, under both Knuth-Bendix ordering and lexicographic path ordering) \citep{...}, the E prover in a proof and a saturation mode \citep{...}, and Twee~2.6.1, an implementation of unfailing completion built for equational reasoning \citep{...}. Each configuration is run over a ladder of per-prover time budgets, $30$, $60$, $120$, $300$, and $600$ seconds, and a law counts as resolved as soon as any configuration returns a verdict. Provers, versions, and exact flags are pinned in a released container image. A claimed model or triviality proof is one the tool itself reports rather than one we infer. We deliberately omit finite-model finders such as Mace4 and Paradox, as admissibility certifies that no nontrivial finite model exists, so these tools cannot succeed. The portfolio's saturation and completion model are instead the automated methods that are able to attempt infinite-model construction. 


% Method-bounded frontier
The hard tier is the set of admissible laws the filtering pass does not resolve. Its difficulty is bound by method rather than computation budget, demonstrated in Table \ref{tab:baseline}. The filtering pass runs a single prover for thirty seconds, building 262 models each in a median of 0.1 seconds. We then reanalyze a sample of the residual with the full eight-configuration portfolio across the budget ladder. If compute were the limit, we would expect resolution to climb as the budget grows. Instead, on 120 laws, the portfolio solves three at the shortest budget, and none thereafter, even at twenty times the time per configuration. 
under every configuration as a ``hard" tier of laws.

The three laws it does settle were all proved trivial and were not yet shown to entail $x = y$ by the short filter. These are contamination that the stronger portfolio removes, not models it builds. The construction side of the hard tier is untouched by the strongest automated methods we can assemble. 

Our solver portfolio is built to contain provers that reason in different ways. Vampire and E are superposition provers: they search by deriving logical consequences of the law until either a contradiction proves it trivial or the consequences close into a model. Twee instead runs unfailing completion, the equational method behind Knuth-Bendix rewriting \citep{...}. Twee turns the law's equations into rewrite rules and completes them into a convergent system. These are two major, distinct paradigms for equational reasoning, each exploring the space differently enough that a problem one diverges on, the other may finish. Twee is also purpose-built for the unit-equational setting our laws fall in. In our experiments, Twee resolves an order-five law in seconds whose saturation runs Vampire to hundreds of clauses. However, across the hard tier, Twee resolves nothing the rest of the portfolio does not, and produces no model of its own. Consistent failures across both automated proof paradigms suggests that the hard frontier's difficulty is shared across automated provers as a class, not particular to any one. 

\input{Figures/table_compute}

\subsection{LLM Baseline}
%LLM Baseline
Todo: Discuss LLM evaluation protocol, show these problems are hard for LLMs too, eval is fair by consrtuction: all LLMs get the same "jumping-off" point to reason from. 

\subsection{Case Study}
Todo:
Our pipeline solves a few open problems mentioned in the ETP paper. We can discuss these two problems, showing that they are confirmed Austin, and how they are solved to give a better understanding of the types of problems we provide. 

% Case Study




\section{Discussion}
% TODO. Anchors: hard-tier collapse not yet measured (prover-only census there);
% ground confluence open => saturation-model answers not yet Lean-certifiable (algebraic
% ones are), framed as next paper; optional hand-solved demonstrations; (i)-prover vs
% Infinox not yet benchmarked.


\bibliography{aaai2027}

% Check whether the conference requires a reproducibility checklist to be included in the paper.
% If so, you can uncomment the following line and ajust the path to include it.
% \input{ReproducibilityChecklist.tex}

\end{document}
